\scanpage{7}
\mspage{L 20}

et plus essentiel. Il était suscité par ces ``interpellations'' par\unsure{f|g}ois déconcertantes, me venant de mon passé \inserted{de mathématicien\deleted{\unsure{s}},} par \deleted{\unsure{la|le} truchement de ceux} voix \deleted{de certains} de ceux qui avaient été mes élèves ou mes amis (ou, du moins, de certains d'entre eux). Au niveau épidermique, ce besoin se traduisait par une envie de \handwritten{``}vider \deleted{un peu} mon sac,\handwritten{''} de dire quelques ``vérités déplaisantes''. Mais plus profondément, sûrement, il y avait le besoin de \mbox{f\,a\,i\,r\,e\quad c\,o\,n\,n\,a\,i\,s\,s\,a\,n\,c\,e} enfin avec un passé, que j'avais choisi jusque là d'éluder. C'est de ce besoin-là, avant tout, qu'est issu Récoltes et Semailles, \deleted{qui}\inserted{lequel} a été ma``réponse'', au jour le jour, à cette pulsion de connaissance en moi,et à l'interpellation sans cesse renouvelée qui me venait du monde extérieur, du``monde \inserted{mathématique''} que j'avais quitté. \deleted{Dès} Mis à part les \handwritten{\inserted{toutes}} premières \deleted{\unsure{deux sept|dix-sept}} pages de \deleted{\unsure{la brève|la partie}} ``Fatuité et Renouvellement'' \inserted{en} formant les deux premiers\deleted{es parties}\inserted{chapitres} (``Travail et découverte'' et ``Le rêve et le Rêveur''), et dès le chapitre qui enchaîne ``Naissance de la crainte'' \deleted{\unsure{par}}\handwritten{\inserted{(p. \unsure{18|10})}}, par lequel commence \deleted{le} un ``témoignage'' qui n'était nullement prévu au programme, c'est ce besoin de faire connaissance de mon passé et de l'assumer pleinement, qui (je crois) a été la force principale en oeuvre dans l'écriture de Récoltes et Semailles.

L'interpellation qui m'était venue du monde des mathématiciens, et qui revenait sur moi avec une force nouvelle tout au cours de l'écriture des parties II,\deleted{\unsure{III,}}\handwritten{\inserted{et}} IV de Récoltes et Semailles, \deleted{prenait} avait pris d'emblée le masque de la suffisance, quand ce n'était celui du dédain (``délicatement dosé''), de la dérision ou du mépris, \handwritten{\inserted{que ce soit}} vis à vis de moi (\deleted{rarement}\handwritten{\inserted{parfois}}) ou \deleted{(à tous les coups ou presque)} \handwritten{\inserted{(surtout)}} vis à vis de ceux qui \deleted{avaient}\inserted{avaient} osé s'inspirer de moi\deleted{\unsure{,}} (sans se douter, \handwritten{\inserted{certes,}} de ce qui les attendait) et qui étaient ``classés'' comme ayant partie liée à moi. Et à nouveau je vois apparaître ici le lien évident et profond, entre le \uline{respect} (ou l'absence \inserted{de respect)}\deleted{\unsure{mépris) de}}\inserted{pour} la personne d'autrui, \deleted{et} celui \deleted{de\unsure{s} des processus créateurs}\inserted{pour l'acte de création} \deleted{primordiaux (dont il vient d'être question) et de leurs fruits, et enfin}\inserted{et pour certains de ses fruits parmi les plus délicats et les plus essentiels, et enfin} le respect pour les règles \deleted{élémentaires d'éthique scientifique}\handwritten{\inserted{les plus évidentes}} de l'éthique scientifique\unsure{,|.} \deleted{Faisant le bilan de ce que m'a \unsure{appris|ap} depuis il y a trois mois, de ce que}\handwritten{\inserted{celles que je serais tenté d'appeler les \uline{``règles de décence''}. Ce sont là autant d'aspects, sûrement, d'un élémentaire et essentiel \uline{respect de soi}.}} Si j'essaie de faire le bilan, en une \deleted{for} seule formule lapidaire, de ce que \unsure{m'a|màa} enseigné\inserted{\unsure{s}} ma réflexion au sujet \deleted{du monde mathéma\unsure{-}} \deleted{\handwritten{\unsure{d'une}}}\deleted{de ce}\handwritten{\inserted{d'un certain}} monde qui fut le mien, auquel je m'étais identifié pendant plus de vingt ans de ma vie, je dirais: c'est un monde qui a \deleted{\uline{perdu le respect}}\mbox{p\,e\,r\,d\,u\quad l\,e\quad r\,e\,s\,p\,e\,c\,t} (\textsuperscript{*}).

% archive sheet number: 6
